% Measured DFC-EF GPU systems insert. Generated from the sealed Kaggle P100
% evidence after docs/DFC_EF_FREE_VALIDATION_PROTOCOL.md was fixed.
% Source kernel commit: edc10ad783aabc2354ef5f1a906634122816d659
% Evidence artifact SHA256:
% acb21bd18be52e7d68b23e49c779ba4d427a5c4b3eeb2490e5f500d9b4192820

\subsection{Actual CUDA Memory Substitution and OOM Frontier}
\label{sec:dfc-ef-gpu}

We next test whether the physical resource substitution predicted by DFC-EF is
visible in an actual GPU allocator rather than only in checkpoint size.  The
criteria were fixed before the accepted run.  The experiment used a Kaggle
Tesla P100-PCIE-16GB (17,059,545,088 reported HBM bytes) with PyTorch
2.7.1+cu126.  The state-only frontier deliberately excludes model weights,
gradients, and activations because these allocations are common to both
placements.  The external baseline owns two FP32 Adam moment containers and one
FP32 error-feedback residual (12 bytes per trainable coordinate), whereas
DFC-EF owns only the two FP32 moment containers (8 bytes per coordinate) and
stores the 32-bit logical residual in their decoder fibers.

\paragraph{CUDA exactness.}
A 1,048,579-coordinate trajectory was executed for 32 updates with a bounded
131,071-coordinate workspace for both FP32 and FP16 model parameters.  The test
includes payload round-trip, decoded optimizer-state equivalence, logical
residual equivalence, parameter equivalence, and checkpoint/restart.  Both
precision paths pass.  During validation we found a real mixed-precision
checkpoint issue: the generic PyTorch optimizer loader casts floating optimizer
state toward the parameter dtype, which is invalid when an FP32 tensor is a
physical DFC container.  The accepted implementation therefore uses a direct
bit-preserving FP32-container reload, covered by a regression test.

\paragraph{Measured $4P$ allocation law.}
Table~\ref{tab:dfc-ef-gpu-allocation} reports CUDA allocator measurements.  At
all five tested scales, the steady allocated-memory difference between external
EF and DFC-EF is exactly $4P$ bytes, while the final parameter and logical
residual digests match.  The smallest case has a transient peak disadvantage
for DFC because the bounded workspace dominates at that scale; at larger scales
the peak allocator difference approaches the same $4P$ law.

\begin{table}[t]
\centering
\caption{Measured P100 allocation substitution.  ``Saved'' is external-EF
minus DFC-EF.  Steady savings equal the predicted $4P$ bytes at every tested
scale.  Timing is retained as a negative result for the unfused reference
implementation.}
\label{tab:dfc-ef-gpu-allocation}
\resizebox{\columnwidth}{!}{%
\begin{tabular}{r|r|r|r|r|r}
\hline
$P$ & steady saved & peak saved & external ms & DFC ms & overhead \\
\hline
1,048,576  & 4,194,304   & -4,194,304  & 0.6813 & 0.9443 & +38.60\% \\
4,194,304  & 16,777,216  & 8,388,608   & 2.2241 & 3.4634 & +55.73\% \\
16,777,216 & 67,108,864  & 58,720,256  & 8.4081 & 13.2998 & +58.18\% \\
33,554,432 & 134,217,728 & 125,829,120 & 16.4633 & 26.5798 & +61.45\% \\
67,108,864 & 268,435,456 & 260,046,848 & 32.9196 & 53.1638 & +61.50\% \\
\hline
\end{tabular}}
\end{table}

\paragraph{OOM frontier.}
Fresh subprocesses were used for every allocation probe so a failed allocation
could not contaminate a later probe.  Only a genuine CUDA OOM return was
classified as OOM; unrelated runtime failures invalidate the probe rather than
counting as evidence.  With 4,194,304-coordinate search resolution, the
external baseline reached 1,371,537,408 successful coordinates, whereas DFC-EF
reached 2,059,403,264, a measured frontier ratio of
\textbf{1.501529$\times$}.  This exceeds the predeclared 1.25$\times$ gate and
is close to the ideal $12/8=1.5$ ratio.

More strongly, at 1,375,731,712 coordinates the external placement requires
16,508,780,544 state bytes and produces an actual CUDA OOM.  At the identical
coordinate count, DFC-EF succeeds with a measured allocation of
11,005,853,696 bytes, exactly matching the predicted $8P$, while providing
5,502,926,848 bytes of hidden residual capacity.  Thus the predeclared
external-OOM/DFC-success crossover gate also passes.  At its maximum measured
frontier, DFC-EF provides 8,237,613,056 bytes of decoder-fiber payload capacity
without a separate model-scale residual allocation.

\paragraph{Performance boundary.}
These measurements establish physical state substitution and a shifted OOM
frontier; they do \emph{not} establish a speed advantage.  The current
bounded-workspace PyTorch reference path is 38.6--61.5\% slower than the matched
external-EF reference over the measured microbenchmark range.  We retain this
negative result explicitly.  Kernel fusion and lower-level implementation are
therefore required before making a deployment-speed claim.  Communication
reduction, separately, is attributable to the matched compressor and not to
DFC itself.
